% s-083-p083-m002; original PDF 83, 84
\Aphorism{s-083-p083-m002}{78}
Die Stärke der „abschliessenden Instincte“ verräth sich in der Ausdehnung der \Sperr{indifferenten} Gebiete. Das moralische Urtheil z. B. ist etwas, wofür die kleinen Leute keine \Foreign{sphincteres} haben; sie können es so wenig halten, wie Kinder den Urin. Bei einem deutschen Philosophen, dem grossen Fürsprecher kleiner Leute, fehlt das moralisch Indifferente gänzlich, wie es in der deutschen Militärmoral fehlt: was nicht befohlen ist, ist verboten. Der religiöse Mensch kann seinen Gott nicht halten; er steckt in Allem, was ihm begegnet, sei es ein Schnupfen oder eine Cursbewegung. Ueberall eine Alternative sehen und das \Foreign{tertium} nicht sehen, ist eine Kinderkrankheit wie das Bettnässen, ein Beweis von \Foreign{impotentia sui}, ein Mangel an Haltung gegenüber zudringlichen Fragestellern und Beichtvätern; man muss sich dieser anmassenden Begriffe und Ideale, die einem das Haus einlaufen und beständig die Rede auf sich bringen, zu erwehren wissen. „Bedaure, der Herr ist nicht zu sprechen“, sollte es immer heissen, wenn die Moral dem erwachsenen Menschen Besuch macht, um ihm „Moral“ zu predigen.
\par

